\section{Dimension and the distribution of the rows of A098568}
\label{pdt:sec:dimension}
Choose a size-$N$ table with probability proportional to $e^{tm}$, and let $D_N$ denote its dimension. Thus
\[
 \Prob_t(D_N=m)=\frac{T(N,m)e^{tm}}{B_N(t)}.
\]
In this section $r,A,\sigma,\lambda_j,C_j$ are evaluated at the same $t$, and all assertions are uniform on each fixed interval $|t|\le T$.

\begin{theorem}[Local expansion and Gaussian limit]\label{pdt:thm:dimension}
Let $\varphi(z)=(2\pi)^{-1/2}e^{-z^2/2}$ and $z=\sqrt A(m-r)$. For every fixed $R<\infty$, uniformly over support points with $|z|\le R$,
\begin{equation}\label{pdt:eq:localexpansion}
 \begin{split}
 \Prob_t(D_N=m)=\sqrt A\,\varphi(z)\bigg[
 1+\frac{\lambda_3z^3}6+\frac{\lambda_4z^4}{24}
 +\frac{\lambda_3^2z^6}{72}-C_1+O_R(N^{-3/2})\bigg].
 \end{split}
\end{equation}
Moreover, with probabilities set to zero off the support,
\begin{equation}\label{pdt:eq:LLT}
 \sup_{m\in\Z}\left|\sigma\Prob_t(D_N=m)
                    -\varphi\!\left(\frac{m-r}{\sigma}\right)\right|\longrightarrow0,
\end{equation}
and
\begin{equation}\label{pdt:eq:CLT}
 \frac{D_N-r}{\sigma}\ \xrightarrow{\ d\ }\ \mathcal N(0,1).
\end{equation}
The support $\{1,\ldots,N\}$ has span one.
\end{theorem}
\begin{proof}
Taylor expansion of the exact numerator at fixed $z$, followed by division by \eqref{pdt:eq:main} through the first correction, gives \eqref{pdt:eq:localexpansion}. The term $-C_1$ is its normalizing correction. The expansion is centered at the saddle, rather than the exact mean, which explains the cubic $z^3$ term.

For $m\in I_N$, \eqref{pdt:eq:gausstail} and the denominator lower bound imply
$\sigma\Prob_t(D_N=m)\le C e^{-cz^2}$. Outside $I_N$, the corresponding bound is $Ce^{-cN}$. The Gaussian itself tends uniformly to zero at those remote support points and outside $\{1,\ldots,N\}$. First apply the compact-$z$ formula and then let $R\to\infty$ to prove \eqref{pdt:eq:LLT}. For the distributional limit, the same tail bound gives tightness, and the compact sums become Gaussian Riemann sums with mesh $1/\sigma\to0$. This proves \eqref{pdt:eq:CLT}; an unweighted supremum local error alone would not suffice. Finally, every $T(N,m)$ is positive, so there is no sublattice obstruction.
\end{proof}

\begin{proposition}[Mean and variance corrections]\label{pdt:prop:moments}
At the same saddle,
\begin{align}
 \E_tD_N&=r+\frac{f_3}{2A^2}+O(\sigma N^{-3/2}),\label{pdt:eq:mean}\\
 \Var_tD_N&=A^{-1}+\frac{f_4}{2A^3}+\frac{f_3^2}{A^4}
                 +O(A^{-1}N^{-2}).\label{pdt:eq:variance}
\end{align}
In particular, at $t=0$,
\[
 \E D_N\sim\frac{2N}{\log N},\qquad
 \Var D_N\sim\frac{2N}{(\log N)^2}.
\]
The error in \eqref{pdt:eq:mean} is $O(1/(N\log N))$.
\end{proposition}
\begin{proof}
Repeat the proof of Section~\ref{pdt:sec:transfer} with factors $z$ and $z^2$. The polynomial tail bounds give uniform integrability, so these weighted expansions are legitimate. Parity cancels the next inappropriate weights and yields
\[
 \E_tz=\lambda_3/2+O(N^{-3/2}),\qquad
 \E_tz^2=1+\lambda_4/2+5\lambda_3^2/4+O(N^{-2}).
\]
For example the unnormalized $z^2$ correction is
$15\lambda_4/24+105\lambda_3^2/72$; subtracting $C_1$ gives the displayed raw moment. Subtracting $(\E_tz)^2$ gives
$\Var_tz=1+\lambda_4/2+\lambda_3^2+O(N^{-2})$.
Rescaling proves both formulas. This proof does not differentiate a real-mark remainder.
\end{proof}
Centering and scaling by the exact mean and standard deviation gives the same normal limit, since their standardized corrections tend to zero. All higher fixed moments and local corrections are obtainable from the weighted-contraction rule and finite series division.

At moderate distances the relative Gaussian approximation holds uniformly whenever $|z|=o(N^{1/6})$. Indeed, these displacements are $o(r)$ and the first nonquadratic exponent is $O(N^{-1/2}|z|^3)=o(1)$ on the intervening proportional interval. Beyond this range it is safer to retain the exact shape: for any support point,
\begin{equation}\label{pdt:eq:exactshape}
 \Prob_t(D_N=m)=\sqrt{\frac A{2\pi}}
 e^{f_N(m)+tm-f_N(r)-tr}
 \left(\sum_{\ell=0}^{K-1}C_\ell+O_{K,T}(N^{-K})\right)^{-1}.
\end{equation}
This is the exact numerator divided by the proved denominator; it does not extend a Gaussian relative approximation into a large-deviation region.

\begin{writenote}[Added 5 October 2026, batch~102]
The moderate-distance claim of the preceding paragraph is argued in two
sentences. The following proposition states it with an explicit error and
proves it from \eqref{pdt:eq:highderivatives}, \eqref{pdt:eq:Ascale} and
\eqref{pdt:eq:exactshape}. The claim ``all higher fixed moments and local
corrections are obtainable'' remains a further question
(Question~\ref{pdt:q:moments}).
\end{writenote}

\begin{proposition}[{[write]} Moderate deviations of the dimension]\label{pdt:prop:moderate}
Fix $T<\infty$. Uniformly for real $|t|\le T$ and for integers $m$ with
$|m-r|\le r/2$, with $z=\sqrt A(m-r)$,
\[
 \Prob_t(D_N=m)=\sqrt A\,\varphi(z)\,
 \exp\!\Bigl(O_T\bigl(|z|^3N^{-1/2}+N^{-1}\bigr)\Bigr).
\]
In particular $\Prob_t(D_N=m)=\sqrt A\,\varphi(z)(1+o(1))$ uniformly over
$|z|\le\omega_NN^{1/6}$, for every sequence $\omega_N\to0$.
\end{proposition}
\begin{proof}
By \eqref{pdt:eq:rscale}, $r=(2N/L)(1+o(1))$ uniformly in $|t|\le T$, so
for all large $N$ every $x$ with $|x-r|\le r/2$ lies in
$I_N=[N/(2L),4N/L]$; in particular such an integer $m$ is a support point.
On $I_N$, \eqref{pdt:eq:highderivatives} with $j=3$ (constants depending
only on the fixed interval) gives $|f_N'''(x)|\le C'N/x^3\le8C'L^3/N^2$,
and \eqref{pdt:eq:Ascale} gives $A\ge L^2/(4N)$, hence
$A^{-3/2}\le8N^{3/2}/L^3$, for all large $N$, uniformly in $t$. Let
$\Phi(x)=f_N(x)+tx$. Since $\Phi'(r)=0$, $\Phi''(r)=-A$ and
$\Phi'''=f_N'''$, Taylor's theorem with the Lagrange remainder gives, for
some $\xi$ between $r$ and $m$,
\[
 \Phi(m)-\Phi(r)=-\frac{z^2}2+\rho,\qquad
 |\rho|=\frac{|f_N'''(\xi)|\,|m-r|^3}6
 \le\frac{8C'L^3}{6N^2}\cdot\frac{8N^{3/2}}{L^3}\,|z|^3
 =\frac{32C'}3\,|z|^3N^{-1/2}.
\]
Formula \eqref{pdt:eq:exactshape} with $K=1$ (so the bracket is
$C_0+O_T(N^{-1})=1+O_T(N^{-1})$) gives
$\Prob_t(D_N=m)=\sqrt{A/(2\pi)}\,e^{-z^2/2+\rho}(1+O_T(N^{-1}))$, which is
the first assertion. For the second, $\sigma=A^{-1/2}\le2N^{1/2}/L$, so
$|z|\le\omega_NN^{1/6}$ implies $|m-r|\le2\omega_NN^{2/3}/L\le r/2$ for
large $N$, while $|z|^3N^{-1/2}\le\omega_N^3\to0$.
\end{proof}
Thus $|z|=o(N^{1/6})$ is the range in which the cubic term
$\lambda_3z^3/6$ of \eqref{pdt:eq:localexpansion}, of size
$\asymp|z|^3N^{-1/2}$, stays small; beyond it the exact shape
\eqref{pdt:eq:exactshape} should be used, as the source says.

To apply the theorem to A098568, normalize its row $n$ by its row sum. If the random column index is $K_n$, then $N=n+1$ and $D_N=K_n+1$. Thus its saddle center is $r_{n+1}(0)-1$, its leading variance is $2(n+1)/(\log(n+1))^2$, and \eqref{pdt:eq:LLT} is a full-span local Gaussian answer to the row-distribution question. The column shift affects the center by one and is not to be dropped in a precise local statement.

\begin{writenote}[Added 5 October 2026, batch~102: the question of A098568]
OEIS A098568 asks, in a comment read on 5~October 2026, ``How do the terms
of row k tend to be distributed as k grows?'' (its row index is our $n$).
Theorem~\ref{pdt:thm:dimension}, Proposition~\ref{pdt:prop:moments} and
Proposition~\ref{pdt:prop:moderate} answer it in the sense above: after
normalization by the row sum $b_{n+1}$, the row is a full-span local
Gaussian law in the column index $K_n=D_{n+1}-1$, with the corrections
\eqref{pdt:eq:localexpansion}, relative accuracy up to $|z|=o(n^{1/6})$ and
the exact shape \eqref{pdt:eq:exactshape} beyond. No repository report
asked this question; it is answered here for the first time in ProveIt, and
the source claims no priority beyond its bounded search. The leading
equivalents are slow: at $n=999$ ($N=1000$) the row has mean column index
$314.606$ and variance $36.557$ (Section~\ref{pdt:sec:checks}), against
$2N/\log N-1=288.5$ and $2N/(\log N)^2=41.9$, while the saddle centre
$r-1=314.514$ and the corrected moments of
Proposition~\ref{pdt:prop:moments} agree with the exact values to about
$7\times10^{-6}$ and $1.5\times10^{-6}$ (intake spot checks, uncertified).
\end{writenote}
